\section{Methodology}
\label{sec:method}

We map person space through four stages: preference item design, persona profile generation, congruence measurement, and dimensionality analysis.

\subsection{Preference Items}
\label{sec:items}

We define \nprefs{} preference items spanning \ndomains{} everyday domains: Food \& Drink, Music \& Arts, Social, Outdoors, Intellectual, Values, Entertainment, Travel, Work, and Technology.
Each domain contributes 6 items (\eg ``loves spicy food,'' ``enjoys abstract art,'' ``thrives in fast-paced environments'').
Every item has a binary polarity: a ``like'' form (\eg ``loves classical music'') and a ``dislike'' form (\eg ``dislikes classical music'').
We deliberately chose everyday preferences rather than psychometric instruments to test whether naturalistic probes reveal structure comparable to formal personality frameworks.

\subsection{Persona Profile Generation}
\label{sec:profiles}

We generate \nprofiles{} persona profiles by randomly selecting 8 of the \nprefs{} preference items, each assigned a random polarity (like or dislike) with equal probability.
The random seed is fixed at 42 for reproducibility.
This combinatorial design ensures rich coverage of the preference space: each pair of items co-occurs in multiple profiles, enabling estimation of pairwise interactions.

\subsection{Congruence Measurement}
\label{sec:congruence}

For each profile, we measure two quantities.

\para{Behavioral congruence.}
We set the model's system prompt to adopt the persona defined by the 8 preferences, then ask 5 follow-up questions designed to elicit preference-relevant behavior (\eg ``What would be your ideal way to spend a free Saturday?'', ``Describe your ideal living environment'').
A separate API call scores each response against the assigned preferences on a 0--10 scale using a structured rubric, with temperature set to 0.0 for deterministic scoring.
The behavioral congruence score for a profile is the average across all 5 questions, normalized to $[0, 1]$.

\para{Coherence rating.}
We ask the model to rate how natural or coherent the preference combination feels on a 0--10 scale, with a brief explanation.
This measures the model's own assessment of whether the preference set constitutes a ``believable person.''
The coherence rating is extracted via structured parsing and normalized to $[0, 1]$.

The combined congruence score averages behavioral congruence and coherence.
We test on two models: \nano (primary, optimized for throughput) and \mini (comparison, stronger model), using temperature 0.7 for persona responses and 0.3 for coherence judgments.

\subsection{Co-occurrence Matrix Construction}
\label{sec:cooccurrence}

We construct a $\nprefs \times \nprefs$ preference co-occurrence matrix $\mC$ where entry $C_{ij}$ is the average congruence score across all profiles in which preferences $i$ and $j$ both appear.
Because each profile contains 8 of \nprefs{} items, each pair co-occurs in approximately $\nprofiles \times \binom{6}{2} / \binom{\nprefs}{2} \approx 5$ profiles on average, providing sufficient coverage for stable estimates.

\subsection{Dimensionality Analysis}
\label{sec:dimanalysis}

We analyze the effective dimensionality of $\mC$ through three methods.

\para{SVD and scree analysis.}
We compute the singular value decomposition $\mC = \mU \mSigma \mV^\top$ and report the number of components needed to explain 80\%, 90\%, and 95\% of the total variance.

\para{Participation ratio.}
We compute the participation ratio $\text{PR} = (\sum_i \sigma_i^2)^2 / \sum_i \sigma_i^4$, which gives an effective dimensionality measure robust to noise~\citep{bose2025lore}.

\para{Random baseline.}
We compare against 1,000 random co-occurrence matrices generated by permuting profile congruence scores, preserving marginal distributions while destroying pairwise structure.
We report the rank and participation ratio of the actual matrix relative to the random distribution.

\subsection{Cross-Model Comparison}
\label{sec:crossmodel}

We compare person-space structure across \nano and \mini using four metrics:
(1)~Pearson correlation between vectorized co-occurrence matrices;
(2)~Mantel permutation test for matrix correlation significance;
(3)~Procrustes analysis aligning the top-$k$ principal components and reporting disparity;
(4)~pairwise cross-correlation of individual principal components to test one-to-one correspondence.
